% ---------------------------------------------------------------
%  شکل: سه سرگروه روش‌شناختی هم‌ترازسازی تصویر
%  چیدمان با مختصات صریح (نه positioning نسبی) تا متن‌ها روی هم نیفتند.
%  نیازمند: tikz + arrows.meta
%  استفاده:  % ---------------------------------------------------------------
%  شکل: سه سرگروه روش‌شناختی هم‌ترازسازی تصویر
%  چیدمان با مختصات صریح (نه positioning نسبی) تا متن‌ها روی هم نیفتند.
%  نیازمند: tikz + arrows.meta
%  استفاده:  % ---------------------------------------------------------------
%  شکل: سه سرگروه روش‌شناختی هم‌ترازسازی تصویر
%  چیدمان با مختصات صریح (نه positioning نسبی) تا متن‌ها روی هم نیفتند.
%  نیازمند: tikz + arrows.meta
%  استفاده:  % ---------------------------------------------------------------
%  شکل: سه سرگروه روش‌شناختی هم‌ترازسازی تصویر
%  چیدمان با مختصات صریح (نه positioning نسبی) تا متن‌ها روی هم نیفتند.
%  نیازمند: tikz + arrows.meta
%  استفاده:  \input{figs/arch_registration}
% ---------------------------------------------------------------
\begingroup
\providecommand{\rsubtext}[1]{{\color{xtire}\tiny #1}}

\tikzset{
  rb/.style   ={draw=xnokhodi, line width=.5pt, fill=white, rounded corners=1.5pt,
                text width=2.55cm, minimum height=1.05cm, align=center,
                inner xsep=1.2mm, inner ysep=1.2mm, font=\scriptsize},
  rnet/.style ={rb, fill=xabikam, draw=xabi},
  rloss/.style={rb, fill=xkamrang},
  rio/.style  ={rb, draw=xkhakestari},
  rflow/.style={->, draw=xtire, line width=.6pt,
                >={Stealth[length=1.7mm,width=1.4mm]}},
  rback/.style={->, draw=xghermez, line width=.5pt, dash pattern=on 1.8pt off 1.4pt,
                >={Stealth[length=1.7mm,width=1.4mm]}},
  rlab/.style ={font=\tiny, color=xkhakestari, inner sep=1.2pt, fill=white},
}

\newcommand{\Rhead}[2]{%
  {\small\textbf{#1}}\par\nobreak\vspace{0.4mm}%
  {\scriptsize\color{xtire}#2}\par\nobreak\vspace{2.2mm}}

\begin{figure}[H]
\centering

% ═════════════════ (الف) کلاسیک شدت‌محور ═════════════════
\Rhead{(الف) خانواده‌ی کلاسیک شدت‌محور \,---\, \lr{elastix} / \lr{SimpleITK}}
      {بدون مرحله‌ی آموزش؛ برای هر جفت تصویر یک بهینه‌سازی تکراری جداگانه اجرا می‌شود}
\begin{tikzpicture}
  \node[rio]   (mov)  at (     0, 0)    {\beginR تصویر متحرک\endR};
  \node[rnet]  (tr)   at (  -3.5, 0)    {\beginR تبدیل $T_\theta$\endR\\[1pt]\rsubtext{\beginR شبکه‌ی نقاط کنترل بی‌اسپلاین\endR}};
  \node[rio]   (warp) at (    -7, 0)    {\beginR تصویر تغییرشکل‌یافته\endR};
  \node[rloss] (sim)  at ( -10.5, 0)    {\beginR معیار شباهت\endR\\[1pt]\rsubtext{\beginR اطلاعات متقابل\endR}};
  \node[rio]   (fix)  at ( -10.5, 1.8)  {\beginR تصویر ثابت\endR};
  \node[rnet]  (opt)  at (    -7,-1.8)  {\beginR بهینه‌ساز گرادیان‌محور\endR};

  \draw[rflow] (mov.west)  -- (tr.east);
  \draw[rflow] (tr.west)   -- (warp.east);
  \draw[rflow] (warp.west) -- (sim.east);
  \draw[rflow] (fix.south) -- (sim.north);
  \draw[rflow] (sim.south) |- (opt.west);
  \draw[rback] (opt.east)  -| node[rlab, pos=0.35] {\beginR به‌روزرسانی $\theta$ و تکرار\endR} (tr.south);
\end{tikzpicture}

\par\vspace{7mm}

% ═════════════════ (ب) توصیف‌گر ساختاری ═════════════════
\Rhead{(ب) توصیف‌گرهای ساختاری \,---\, \lr{MIND} / \lr{DEEDS}}
      {به‌جای \emph{مقدار} شدت روشنایی، \emph{الگوی محلی} آن مقایسه می‌شود}
\begin{tikzpicture}
  \node[rio]   (mv)  at (    0, 1.0)  {\beginR تصویر متحرک\endR};
  \node[rio]   (fx)  at (    0,-1.0)  {\beginR تصویر ثابت\endR};
  \node[rnet]  (d1)  at ( -3.5, 1.0)  {\beginR توصیف‌گر \lr{MIND}\endR\\[1pt]\rsubtext{\beginR خودشباهت وصله‌ها\endR}};
  \node[rnet]  (d2)  at ( -3.5,-1.0)  {\beginR توصیف‌گر \lr{MIND}\endR};
  \node[rloss] (ssd) at (   -7, 0)    {\beginR \lr{SSD} روی بردار توصیف‌گر\endR};
  \node[rnet]  (mrf) at (-10.5, 0)    {\beginR بهینه‌سازی گسسته‌ی \lr{MRF}\endR\\[1pt]\rsubtext{\beginR انتقال پیام روی گراف\endR}};
  \node[rio]   (phi) at (-10.5,-2.0)  {\beginR میدان جابه‌جایی\endR};

  \draw[rflow] (mv.west) -- (d1.east);
  \draw[rflow] (fx.west) -- (d2.east);
  \draw[rflow] (d1.west) -- ++(-0.55,0) |- (ssd.east);
  \draw[rflow] (d2.west) -- ++(-0.55,0) |- (ssd.east);
  \draw[rflow] (ssd.west)  -- (mrf.east);
  \draw[rflow] (mrf.south) -- (phi.north);
\end{tikzpicture}

\par\vspace{7mm}

% ═════════════════ (پ) یادگیری‌محور ═════════════════
\Rhead{(پ) خانواده‌ی یادگیری‌محور \,---\, \lr{VoxelMorph} / \lr{TransMorph}}
      {آموزش یک‌بار روی کل مجموعه؛ استنتاج در یک گذر پیش‌رو}
\begin{tikzpicture}
  \node[rio]   (in)   at (     0, 0)    {\beginR متحرک $\oplus$ ثابت\endR\\[1pt]\rsubtext{\beginR الحاق در بعد کانال\endR}};
  \node[rnet]  (enc)  at (  -3.5, 0)    {\beginR رمزگذار--رمزگشا\endR\\[1pt]\rsubtext{\beginR \lr{U-Net} یا \lr{Swin}\endR}};
  \node[rio]   (phi3) at (    -7, 0)    {\beginR میدان جابه‌جایی $\phi$\endR};
  \node[rnet]  (stn)  at ( -10.5, 0)    {\beginR تبدیل‌گر فضایی\endR\\[1pt]\rsubtext{مشتق‌پذیر}};
  \node[rloss] (ls)   at (    -7,-1.9)  {\beginR زیان: شباهت تصویر $+$ هموارسازی $\phi$\endR};

  \draw[rflow] (in.west)   -- (enc.east);
  \draw[rflow] (enc.west)  -- (phi3.east);
  \draw[rflow] (phi3.west) -- (stn.east);
  \draw[rflow] (stn.south) |- (ls.west);
  \draw[rback] (ls.east)   -| node[rlab, pos=0.35] {\beginR پس‌انتشار (فقط در آموزش)\endR} (enc.south);
\end{tikzpicture}

\caption{سه سرگروه روش‌شناختی هم‌ترازسازی تصویر: (الف) معیارهای شدت‌محور، (ب) توصیف‌گرهای ساختاری، (پ) روش‌های یادگیری‌محور.}
\label{fig:reg-archetypes}
\end{figure}
\endgroup

% ---------------------------------------------------------------
\begingroup
\providecommand{\rsubtext}[1]{{\color{xtire}\tiny #1}}

\tikzset{
  rb/.style   ={draw=xnokhodi, line width=.5pt, fill=white, rounded corners=1.5pt,
                text width=2.55cm, minimum height=1.05cm, align=center,
                inner xsep=1.2mm, inner ysep=1.2mm, font=\scriptsize},
  rnet/.style ={rb, fill=xabikam, draw=xabi},
  rloss/.style={rb, fill=xkamrang},
  rio/.style  ={rb, draw=xkhakestari},
  rflow/.style={->, draw=xtire, line width=.6pt,
                >={Stealth[length=1.7mm,width=1.4mm]}},
  rback/.style={->, draw=xghermez, line width=.5pt, dash pattern=on 1.8pt off 1.4pt,
                >={Stealth[length=1.7mm,width=1.4mm]}},
  rlab/.style ={font=\tiny, color=xkhakestari, inner sep=1.2pt, fill=white},
}

\newcommand{\Rhead}[2]{%
  {\small\textbf{#1}}\par\nobreak\vspace{0.4mm}%
  {\scriptsize\color{xtire}#2}\par\nobreak\vspace{2.2mm}}

\begin{figure}[H]
\centering

% ═════════════════ (الف) کلاسیک شدت‌محور ═════════════════
\Rhead{(الف) خانواده‌ی کلاسیک شدت‌محور \,---\, \lr{elastix} / \lr{SimpleITK}}
      {بدون مرحله‌ی آموزش؛ برای هر جفت تصویر یک بهینه‌سازی تکراری جداگانه اجرا می‌شود}
\begin{tikzpicture}
  \node[rio]   (mov)  at (     0, 0)    {\beginR تصویر متحرک\endR};
  \node[rnet]  (tr)   at (  -3.5, 0)    {\beginR تبدیل $T_\theta$\endR\\[1pt]\rsubtext{\beginR شبکه‌ی نقاط کنترل بی‌اسپلاین\endR}};
  \node[rio]   (warp) at (    -7, 0)    {\beginR تصویر تغییرشکل‌یافته\endR};
  \node[rloss] (sim)  at ( -10.5, 0)    {\beginR معیار شباهت\endR\\[1pt]\rsubtext{\beginR اطلاعات متقابل\endR}};
  \node[rio]   (fix)  at ( -10.5, 1.8)  {\beginR تصویر ثابت\endR};
  \node[rnet]  (opt)  at (    -7,-1.8)  {\beginR بهینه‌ساز گرادیان‌محور\endR};

  \draw[rflow] (mov.west)  -- (tr.east);
  \draw[rflow] (tr.west)   -- (warp.east);
  \draw[rflow] (warp.west) -- (sim.east);
  \draw[rflow] (fix.south) -- (sim.north);
  \draw[rflow] (sim.south) |- (opt.west);
  \draw[rback] (opt.east)  -| node[rlab, pos=0.35] {\beginR به‌روزرسانی $\theta$ و تکرار\endR} (tr.south);
\end{tikzpicture}

\par\vspace{7mm}

% ═════════════════ (ب) توصیف‌گر ساختاری ═════════════════
\Rhead{(ب) توصیف‌گرهای ساختاری \,---\, \lr{MIND} / \lr{DEEDS}}
      {به‌جای \emph{مقدار} شدت روشنایی، \emph{الگوی محلی} آن مقایسه می‌شود}
\begin{tikzpicture}
  \node[rio]   (mv)  at (    0, 1.0)  {\beginR تصویر متحرک\endR};
  \node[rio]   (fx)  at (    0,-1.0)  {\beginR تصویر ثابت\endR};
  \node[rnet]  (d1)  at ( -3.5, 1.0)  {\beginR توصیف‌گر \lr{MIND}\endR\\[1pt]\rsubtext{\beginR خودشباهت وصله‌ها\endR}};
  \node[rnet]  (d2)  at ( -3.5,-1.0)  {\beginR توصیف‌گر \lr{MIND}\endR};
  \node[rloss] (ssd) at (   -7, 0)    {\beginR \lr{SSD} روی بردار توصیف‌گر\endR};
  \node[rnet]  (mrf) at (-10.5, 0)    {\beginR بهینه‌سازی گسسته‌ی \lr{MRF}\endR\\[1pt]\rsubtext{\beginR انتقال پیام روی گراف\endR}};
  \node[rio]   (phi) at (-10.5,-2.0)  {\beginR میدان جابه‌جایی\endR};

  \draw[rflow] (mv.west) -- (d1.east);
  \draw[rflow] (fx.west) -- (d2.east);
  \draw[rflow] (d1.west) -- ++(-0.55,0) |- (ssd.east);
  \draw[rflow] (d2.west) -- ++(-0.55,0) |- (ssd.east);
  \draw[rflow] (ssd.west)  -- (mrf.east);
  \draw[rflow] (mrf.south) -- (phi.north);
\end{tikzpicture}

\par\vspace{7mm}

% ═════════════════ (پ) یادگیری‌محور ═════════════════
\Rhead{(پ) خانواده‌ی یادگیری‌محور \,---\, \lr{VoxelMorph} / \lr{TransMorph}}
      {آموزش یک‌بار روی کل مجموعه؛ استنتاج در یک گذر پیش‌رو}
\begin{tikzpicture}
  \node[rio]   (in)   at (     0, 0)    {\beginR متحرک $\oplus$ ثابت\endR\\[1pt]\rsubtext{\beginR الحاق در بعد کانال\endR}};
  \node[rnet]  (enc)  at (  -3.5, 0)    {\beginR رمزگذار--رمزگشا\endR\\[1pt]\rsubtext{\beginR \lr{U-Net} یا \lr{Swin}\endR}};
  \node[rio]   (phi3) at (    -7, 0)    {\beginR میدان جابه‌جایی $\phi$\endR};
  \node[rnet]  (stn)  at ( -10.5, 0)    {\beginR تبدیل‌گر فضایی\endR\\[1pt]\rsubtext{مشتق‌پذیر}};
  \node[rloss] (ls)   at (    -7,-1.9)  {\beginR زیان: شباهت تصویر $+$ هموارسازی $\phi$\endR};

  \draw[rflow] (in.west)   -- (enc.east);
  \draw[rflow] (enc.west)  -- (phi3.east);
  \draw[rflow] (phi3.west) -- (stn.east);
  \draw[rflow] (stn.south) |- (ls.west);
  \draw[rback] (ls.east)   -| node[rlab, pos=0.35] {\beginR پس‌انتشار (فقط در آموزش)\endR} (enc.south);
\end{tikzpicture}

\caption{سه سرگروه روش‌شناختی هم‌ترازسازی تصویر: (الف) معیارهای شدت‌محور، (ب) توصیف‌گرهای ساختاری، (پ) روش‌های یادگیری‌محور.}
\label{fig:reg-archetypes}
\end{figure}
\endgroup

% ---------------------------------------------------------------
\begingroup
\providecommand{\rsubtext}[1]{{\color{xtire}\tiny #1}}

\tikzset{
  rb/.style   ={draw=xnokhodi, line width=.5pt, fill=white, rounded corners=1.5pt,
                text width=2.55cm, minimum height=1.05cm, align=center,
                inner xsep=1.2mm, inner ysep=1.2mm, font=\scriptsize},
  rnet/.style ={rb, fill=xabikam, draw=xabi},
  rloss/.style={rb, fill=xkamrang},
  rio/.style  ={rb, draw=xkhakestari},
  rflow/.style={->, draw=xtire, line width=.6pt,
                >={Stealth[length=1.7mm,width=1.4mm]}},
  rback/.style={->, draw=xghermez, line width=.5pt, dash pattern=on 1.8pt off 1.4pt,
                >={Stealth[length=1.7mm,width=1.4mm]}},
  rlab/.style ={font=\tiny, color=xkhakestari, inner sep=1.2pt, fill=white},
}

\newcommand{\Rhead}[2]{%
  {\small\textbf{#1}}\par\nobreak\vspace{0.4mm}%
  {\scriptsize\color{xtire}#2}\par\nobreak\vspace{2.2mm}}

\begin{figure}[H]
\centering

% ═════════════════ (الف) کلاسیک شدت‌محور ═════════════════
\Rhead{(الف) خانواده‌ی کلاسیک شدت‌محور \,---\, \lr{elastix} / \lr{SimpleITK}}
      {بدون مرحله‌ی آموزش؛ برای هر جفت تصویر یک بهینه‌سازی تکراری جداگانه اجرا می‌شود}
\begin{tikzpicture}
  \node[rio]   (mov)  at (     0, 0)    {\beginR تصویر متحرک\endR};
  \node[rnet]  (tr)   at (  -3.5, 0)    {\beginR تبدیل $T_\theta$\endR\\[1pt]\rsubtext{\beginR شبکه‌ی نقاط کنترل بی‌اسپلاین\endR}};
  \node[rio]   (warp) at (    -7, 0)    {\beginR تصویر تغییرشکل‌یافته\endR};
  \node[rloss] (sim)  at ( -10.5, 0)    {\beginR معیار شباهت\endR\\[1pt]\rsubtext{\beginR اطلاعات متقابل\endR}};
  \node[rio]   (fix)  at ( -10.5, 1.8)  {\beginR تصویر ثابت\endR};
  \node[rnet]  (opt)  at (    -7,-1.8)  {\beginR بهینه‌ساز گرادیان‌محور\endR};

  \draw[rflow] (mov.west)  -- (tr.east);
  \draw[rflow] (tr.west)   -- (warp.east);
  \draw[rflow] (warp.west) -- (sim.east);
  \draw[rflow] (fix.south) -- (sim.north);
  \draw[rflow] (sim.south) |- (opt.west);
  \draw[rback] (opt.east)  -| node[rlab, pos=0.35] {\beginR به‌روزرسانی $\theta$ و تکرار\endR} (tr.south);
\end{tikzpicture}

\par\vspace{7mm}

% ═════════════════ (ب) توصیف‌گر ساختاری ═════════════════
\Rhead{(ب) توصیف‌گرهای ساختاری \,---\, \lr{MIND} / \lr{DEEDS}}
      {به‌جای \emph{مقدار} شدت روشنایی، \emph{الگوی محلی} آن مقایسه می‌شود}
\begin{tikzpicture}
  \node[rio]   (mv)  at (    0, 1.0)  {\beginR تصویر متحرک\endR};
  \node[rio]   (fx)  at (    0,-1.0)  {\beginR تصویر ثابت\endR};
  \node[rnet]  (d1)  at ( -3.5, 1.0)  {\beginR توصیف‌گر \lr{MIND}\endR\\[1pt]\rsubtext{\beginR خودشباهت وصله‌ها\endR}};
  \node[rnet]  (d2)  at ( -3.5,-1.0)  {\beginR توصیف‌گر \lr{MIND}\endR};
  \node[rloss] (ssd) at (   -7, 0)    {\beginR \lr{SSD} روی بردار توصیف‌گر\endR};
  \node[rnet]  (mrf) at (-10.5, 0)    {\beginR بهینه‌سازی گسسته‌ی \lr{MRF}\endR\\[1pt]\rsubtext{\beginR انتقال پیام روی گراف\endR}};
  \node[rio]   (phi) at (-10.5,-2.0)  {\beginR میدان جابه‌جایی\endR};

  \draw[rflow] (mv.west) -- (d1.east);
  \draw[rflow] (fx.west) -- (d2.east);
  \draw[rflow] (d1.west) -- ++(-0.55,0) |- (ssd.east);
  \draw[rflow] (d2.west) -- ++(-0.55,0) |- (ssd.east);
  \draw[rflow] (ssd.west)  -- (mrf.east);
  \draw[rflow] (mrf.south) -- (phi.north);
\end{tikzpicture}

\par\vspace{7mm}

% ═════════════════ (پ) یادگیری‌محور ═════════════════
\Rhead{(پ) خانواده‌ی یادگیری‌محور \,---\, \lr{VoxelMorph} / \lr{TransMorph}}
      {آموزش یک‌بار روی کل مجموعه؛ استنتاج در یک گذر پیش‌رو}
\begin{tikzpicture}
  \node[rio]   (in)   at (     0, 0)    {\beginR متحرک $\oplus$ ثابت\endR\\[1pt]\rsubtext{\beginR الحاق در بعد کانال\endR}};
  \node[rnet]  (enc)  at (  -3.5, 0)    {\beginR رمزگذار--رمزگشا\endR\\[1pt]\rsubtext{\beginR \lr{U-Net} یا \lr{Swin}\endR}};
  \node[rio]   (phi3) at (    -7, 0)    {\beginR میدان جابه‌جایی $\phi$\endR};
  \node[rnet]  (stn)  at ( -10.5, 0)    {\beginR تبدیل‌گر فضایی\endR\\[1pt]\rsubtext{مشتق‌پذیر}};
  \node[rloss] (ls)   at (    -7,-1.9)  {\beginR زیان: شباهت تصویر $+$ هموارسازی $\phi$\endR};

  \draw[rflow] (in.west)   -- (enc.east);
  \draw[rflow] (enc.west)  -- (phi3.east);
  \draw[rflow] (phi3.west) -- (stn.east);
  \draw[rflow] (stn.south) |- (ls.west);
  \draw[rback] (ls.east)   -| node[rlab, pos=0.35] {\beginR پس‌انتشار (فقط در آموزش)\endR} (enc.south);
\end{tikzpicture}

\caption{سه سرگروه روش‌شناختی هم‌ترازسازی تصویر: (الف) معیارهای شدت‌محور، (ب) توصیف‌گرهای ساختاری، (پ) روش‌های یادگیری‌محور.}
\label{fig:reg-archetypes}
\end{figure}
\endgroup

% ---------------------------------------------------------------
\begingroup
\providecommand{\rsubtext}[1]{{\color{xtire}\tiny #1}}

\tikzset{
  rb/.style   ={draw=xnokhodi, line width=.5pt, fill=white, rounded corners=1.5pt,
                text width=2.55cm, minimum height=1.05cm, align=center,
                inner xsep=1.2mm, inner ysep=1.2mm, font=\scriptsize},
  rnet/.style ={rb, fill=xabikam, draw=xabi},
  rloss/.style={rb, fill=xkamrang},
  rio/.style  ={rb, draw=xkhakestari},
  rflow/.style={->, draw=xtire, line width=.6pt,
                >={Stealth[length=1.7mm,width=1.4mm]}},
  rback/.style={->, draw=xghermez, line width=.5pt, dash pattern=on 1.8pt off 1.4pt,
                >={Stealth[length=1.7mm,width=1.4mm]}},
  rlab/.style ={font=\tiny, color=xkhakestari, inner sep=1.2pt, fill=white},
}

\newcommand{\Rhead}[2]{%
  {\small\textbf{#1}}\par\nobreak\vspace{0.4mm}%
  {\scriptsize\color{xtire}#2}\par\nobreak\vspace{2.2mm}}

\begin{figure}[H]
\centering

% ═════════════════ (الف) کلاسیک شدت‌محور ═════════════════
\Rhead{(الف) خانواده‌ی کلاسیک شدت‌محور \,---\, \lr{elastix} / \lr{SimpleITK}}
      {بدون مرحله‌ی آموزش؛ برای هر جفت تصویر یک بهینه‌سازی تکراری جداگانه اجرا می‌شود}
\begin{tikzpicture}
  \node[rio]   (mov)  at (     0, 0)    {\beginR تصویر متحرک\endR};
  \node[rnet]  (tr)   at (  -3.5, 0)    {\beginR تبدیل $T_\theta$\endR\\[1pt]\rsubtext{\beginR شبکه‌ی نقاط کنترل بی‌اسپلاین\endR}};
  \node[rio]   (warp) at (    -7, 0)    {\beginR تصویر تغییرشکل‌یافته\endR};
  \node[rloss] (sim)  at ( -10.5, 0)    {\beginR معیار شباهت\endR\\[1pt]\rsubtext{\beginR اطلاعات متقابل\endR}};
  \node[rio]   (fix)  at ( -10.5, 1.8)  {\beginR تصویر ثابت\endR};
  \node[rnet]  (opt)  at (    -7,-1.8)  {\beginR بهینه‌ساز گرادیان‌محور\endR};

  \draw[rflow] (mov.west)  -- (tr.east);
  \draw[rflow] (tr.west)   -- (warp.east);
  \draw[rflow] (warp.west) -- (sim.east);
  \draw[rflow] (fix.south) -- (sim.north);
  \draw[rflow] (sim.south) |- (opt.west);
  \draw[rback] (opt.east)  -| node[rlab, pos=0.35] {\beginR به‌روزرسانی $\theta$ و تکرار\endR} (tr.south);
\end{tikzpicture}

\par\vspace{7mm}

% ═════════════════ (ب) توصیف‌گر ساختاری ═════════════════
\Rhead{(ب) توصیف‌گرهای ساختاری \,---\, \lr{MIND} / \lr{DEEDS}}
      {به‌جای \emph{مقدار} شدت روشنایی، \emph{الگوی محلی} آن مقایسه می‌شود}
\begin{tikzpicture}
  \node[rio]   (mv)  at (    0, 1.0)  {\beginR تصویر متحرک\endR};
  \node[rio]   (fx)  at (    0,-1.0)  {\beginR تصویر ثابت\endR};
  \node[rnet]  (d1)  at ( -3.5, 1.0)  {\beginR توصیف‌گر \lr{MIND}\endR\\[1pt]\rsubtext{\beginR خودشباهت وصله‌ها\endR}};
  \node[rnet]  (d2)  at ( -3.5,-1.0)  {\beginR توصیف‌گر \lr{MIND}\endR};
  \node[rloss] (ssd) at (   -7, 0)    {\beginR \lr{SSD} روی بردار توصیف‌گر\endR};
  \node[rnet]  (mrf) at (-10.5, 0)    {\beginR بهینه‌سازی گسسته‌ی \lr{MRF}\endR\\[1pt]\rsubtext{\beginR انتقال پیام روی گراف\endR}};
  \node[rio]   (phi) at (-10.5,-2.0)  {\beginR میدان جابه‌جایی\endR};

  \draw[rflow] (mv.west) -- (d1.east);
  \draw[rflow] (fx.west) -- (d2.east);
  \draw[rflow] (d1.west) -- ++(-0.55,0) |- (ssd.east);
  \draw[rflow] (d2.west) -- ++(-0.55,0) |- (ssd.east);
  \draw[rflow] (ssd.west)  -- (mrf.east);
  \draw[rflow] (mrf.south) -- (phi.north);
\end{tikzpicture}

\par\vspace{7mm}

% ═════════════════ (پ) یادگیری‌محور ═════════════════
\Rhead{(پ) خانواده‌ی یادگیری‌محور \,---\, \lr{VoxelMorph} / \lr{TransMorph}}
      {آموزش یک‌بار روی کل مجموعه؛ استنتاج در یک گذر پیش‌رو}
\begin{tikzpicture}
  \node[rio]   (in)   at (     0, 0)    {\beginR متحرک $\oplus$ ثابت\endR\\[1pt]\rsubtext{\beginR الحاق در بعد کانال\endR}};
  \node[rnet]  (enc)  at (  -3.5, 0)    {\beginR رمزگذار--رمزگشا\endR\\[1pt]\rsubtext{\beginR \lr{U-Net} یا \lr{Swin}\endR}};
  \node[rio]   (phi3) at (    -7, 0)    {\beginR میدان جابه‌جایی $\phi$\endR};
  \node[rnet]  (stn)  at ( -10.5, 0)    {\beginR تبدیل‌گر فضایی\endR\\[1pt]\rsubtext{مشتق‌پذیر}};
  \node[rloss] (ls)   at (    -7,-1.9)  {\beginR زیان: شباهت تصویر $+$ هموارسازی $\phi$\endR};

  \draw[rflow] (in.west)   -- (enc.east);
  \draw[rflow] (enc.west)  -- (phi3.east);
  \draw[rflow] (phi3.west) -- (stn.east);
  \draw[rflow] (stn.south) |- (ls.west);
  \draw[rback] (ls.east)   -| node[rlab, pos=0.35] {\beginR پس‌انتشار (فقط در آموزش)\endR} (enc.south);
\end{tikzpicture}

\caption{سه سرگروه روش‌شناختی هم‌ترازسازی تصویر: (الف) معیارهای شدت‌محور، (ب) توصیف‌گرهای ساختاری، (پ) روش‌های یادگیری‌محور.}
\label{fig:reg-archetypes}
\end{figure}
\endgroup
